\section{Synthesis methods}
\label{sec:synthesis}
\textbf{Top-down} methods break bulk material into small pieces: simple, but poor size control. \textbf{Bottom-up} methods build
particles from atoms or molecules (atoms $\to$ clusters $\to$ particles): more complex, but good size and shape control
\slide{29}. In solution, size is set by the competition between \textbf{nucleation and growth}: fast, homogeneous nucleation
(strong reductant, good stirring) gives many small particles; capping agents stop growth; confinement (micelle, pore, droplet)
limits the final size.

Each card lists: \textbf{what it is}, \textbf{advantages} (+) and \textbf{drawbacks} (--), and \textbf{use for} (typical products
and exercises).

\newcommand{\meth}[5]{%
\begin{cardbox}{#1}
\textbf{What it is:} #2\par\smallskip
\textcolor{proc}{\textbf{+}} #3\par
\textcolor{conc}{\textbf{--}} #4\par
\textcolor{recc}{\textbf{Use for:}} #5
\end{cardbox}}

\subsection{Top-down}
\meth{Ball milling \slide{31--32}}
{Powder is sealed in a jar with hard balls and rotated at high speed (planetary mill). The repeated impacts deform the grains,
first elastically, then plastically and by shear, until they fracture. Long milling can also amorphise the material or drive
solid-state reactions (mechanochemistry), e.g.\ mixing CuO with CeO$_2$.}
{cheap, fast, easy, no solvents, scalable.}
{poor size control, contamination from balls/jar, many defects (dislocations).}
{large amounts of powders and mixed oxides (alternative route in Ex.~10e).}
\meth{Lithography + selective etching \slide{33--36}}
{A Si wafer is masked by lithography and etched (TMAH, then photo-assisted etching) into a regular array of
$\sim$\SI{4e6}{channels/cm^2}. The walls are oxidised to SiO$_2$; their --OH groups react with a metal precursor (grafting),
which is calcined to a thin oxide layer that can carry metal nanoparticles (Au/TiO$_2$, RhPd/CeO$_2$).}
{perfect, ordered geometry; very reproducible.}
{multi-step, expensive (clean room).}
{structured supports and micro-monolith catalysts.}

\subsection{Bottom-up, gas phase}
\textbf{PVD} (physical vapour deposition) involves no catalytic interaction; \textbf{CVD} (chemical vapour deposition) involves
a catalytic reaction of the precursor on a surface \slide{38}. The common problem of gas-phase methods is coalescence and
agglomeration of the particles in flight, solved with a carrier gas and fast cooling.

\meth{Sputtering (PVD) \slide{39}}
{An Ar plasma is created in a vacuum chamber. Ar$^+$ ions are accelerated onto a target of the material; momentum transfer ejects
atoms or clusters, which land on a substrate. Short deposition gives isolated particles; longer deposition gives islands and
then a continuous film.}
{clean (no solvents or ligands), works for any metal or alloy target.}
{vacuum equipment, little shape control, small amounts.}
{model catalysts and thin films.}
\meth{Spray pyrolysis \slide{40--46}}
{A solution of a volatile precursor (TiCl$_4$, SiH$_4$) is atomised into droplets and carried by a gas (Ar, N$_2$, O$_2$, H$_2$)
through a furnace at \SIrange{300}{1200}{\celsius}. Each droplet evaporates, precipitates, dries, decomposes and sinters
into one particle (\ce{TiCl4 + O2 -> TiO2 + 2Cl2}); particles are cooled in inert gas. Size is set by temperature, pressure,
gas flow (= reaction time) and concentration: ``the growth is not predictable, but it is reproducible''. Adding polystyrene spheres
that burn off at \SI{600}{\celsius} gives ordered macroporous particles \slide{45}.}
{continuous and scalable; the droplet composition fixes the particle composition (mixed oxides); templates give ordered porosity.}
{agglomeration, empirical tuning, high temperature.}
{oxide particles, porous and photonic structures (Ex.~11c).}
\meth{Flame synthesis \slide{47--48}}
{The precursor is injected into a CH$_4$/O$_2$, H$_2$/O$_2$ or C$_2$H$_2$/O$_2$ flame at \SIrange{1000}{2400}{\celsius}. It reacts into
monomers, which grow by coagulation and coalescence during a residence time of only \SIrange{10}{100}{ms}, and agglomerate after the
flame.}
{one-step, the largest production scale ($10^6$ t/year of SiO$_2$, Al$_2$O$_3$, TiO$_2$; Fe$_3$O$_4$ for ferrofluids).}
{needs fuel; particles form fractal agglomerates; little shape control.}
{bulk oxide powders.}
\meth{Laser processes \slide{49--51}}
{A pulsed laser (\SIrange{10}{50}{ns}, e.g.\ a CO$_2$ laser) crosses a stream of reactant gas or vaporises a target. Heating is
confined to a tiny spot and followed by very fast cooling, so nucleation is short and controlled. With a mass spectrometer one can
select clusters of an exact size (M$_{13}$, M$_{55}$).}
{excellent control of size and morphology, highly localised.}
{expensive energy, small scale (\SIrange{10}{100}{mg/h}).}
{model clusters, SiGe particles.}
\meth{Thermal plasma \slide{52}}
{An arc or magnetron plasma (\SIrange{10}{100}{V}, \SIrange{10}{1000}{A/cm^2}) vaporises the material at $\sim$\SI{e4}{\celsius} for
less than 1 s; the vapour is quenched at \SIrange{e6}{e9}{\celsius/s} in an aggregation zone and leaves as a nanoparticle beam.}
{size control from 1 to 100 nm; refractory and metastable phases.}
{high energy cost, complex equipment.}
{refractory or metastable particles.}
\meth{Atomic layer deposition (ALD) \slide{53--54}}
{Two precursors are pulsed alternately onto the substrate. Each reacts only with the surface groups left by the other and stops when
the surface is saturated (\textbf{self-limiting}), so exactly one atomic layer grows per cycle. Example: trimethylaluminium reacts with
surface --OH (releasing CH$_4$), then H$_2$O turns the --CH$_3$ back into --OH, giving Al$_2$O$_3$ at \SIrange{100}{300}{\celsius}.
Precursors must be volatile, react with the surface, not with themselves, be thermally stable and give volatile by-products.}
{atomic control of thickness and composition; \textbf{conformal} coatings on complex 3D shapes (pores, channels).}
{slow (one layer per cycle), vacuum.}
{ultrathin shells and coatings.}
\meth{CVD and carbon nanotubes \slide{104--121}}
{A hydrocarbon (with an inert gas so it does not oxidise) flows over metal catalyst particles (Ni, Co, Fe) at
\SIrange{500}{1000}{\celsius}. Carbon dissolves in the particle (forming a carbide) and precipitates as graphitic walls. The catalyst
size fixes the tube: 1--5 nm $\to$ single-walled (SWNT), 20--50 nm $\to$ multi-walled (MWNT), $>$50 nm $\to$ nano-onions. A weak
metal--support interaction lifts the particle (tip growth); a strong one keeps it on the support (base growth). The roll-up
vector $(n,m)$ (zig-zag, armchair, chiral) sets the electronic, optical and mechanical properties.}
{yield $\approx$97\,\%, long and aligned tubes, diameter controlled by the catalyst.}
{needs a catalyst. Alternatives: arc discharge (simplest, SWNT + MWNT, no size control, 30\,\%) and laser vaporisation
(high-purity SWNT, narrow distribution, low yield).}
{carbon nanotubes and films.}

\subsection{Bottom-up, liquid phase}
\meth{Reduction of metal salts \slide{56--58}}
{Metal ions in solution are reduced to neutral atoms (\ce{M^{n+} + n e^- -> M^0}) by citrate, borohydride or alcohols. The atoms
nucleate and grow into a colloid. \textbf{Turkevich method}: HAuCl$_4$ is boiled and sodium citrate added; the solution turns
from pale purple to red and gives Au 15--20 nm. Citrate is both reductant and capping agent: more citrate gives smaller particles
(8 nm, 518 nm band $\to$ 73 nm, 545 nm band) \slide{128}. Ag colloids (20--100 nm) are made the same way from AgNO$_3$.}
{simple, reproducible, cheap, stable colloids; size tunable through the reductant/metal ratio, $T$ and stirring.}
{mainly noble metals; mostly spherical particles.}
{Au, Ag and Pt colloids; cores of core--shell particles (Ex.~11b, 13, 15a).}
\meth{Bimetallic nanoparticles \slide{59}}
{\textbf{Co-reduction} of two salts at the same time gives an \textbf{alloy}. \textbf{Sequential (seeded) reduction}, i.e.\ making
M$_1$ first and then reducing M$_2$ on it, gives a \textbf{core--shell} particle (M$_1$@M$_2$). Two segregated halves form a
\textbf{Janus} particle.}
{control of composition and architecture.}
{several techniques are needed to prove which architecture was obtained.}
{Ex.~1, 9, 16b.}
\meth{Hydrothermal / solvothermal \slide{60--61}}
{The reaction runs in a closed autoclave with the solvent (water or an organic solvent) above its boiling point. The solvent
provides the reaction medium and the pressure, and its conditions control size, shape and crystallinity. For CeO$_2$, high NaOH and
high $T$ ($\sim$\SI{180}{\celsius}) give \textbf{nanocubes}, intermediate conditions give \textbf{nanorods}, and mild conditions
give small polyhedra.}
{shape control without organic ligands; crystalline particles with clean facets.}
{batch process, long times, pressure equipment.}
{shaped oxides: CeO$_2$ cubes/rods, CoMg cubes (Ex.~11a, 12A).}
\meth{Sol--gel \slide{62--67}}
{A metal alkoxide such as Si(OR)$_4$ is \textbf{hydrolysed} (\ce{-M-OR + H2O -> -M-OH + ROH}) and \textbf{condensed}
(\ce{-M-OH + HO-M- -> -M-O-M- + H2O}) into a colloidal \textbf{sol}, which links into a 3D \textbf{gel}. Acid conditions give
chain-like networks; basic conditions give discrete particles of 1--100 nm. Catalytic particles can be added to the sol.
The drying step decides the porosity: evaporation (surface tension collapses the pores) gives a \textbf{xerogel}; freeze-drying
(sublimation) gives a \textbf{liogel}; \textbf{supercritical drying} (no liquid/gas interface) gives an \textbf{aerogel}, with
$S_\mathrm{BET}$(xerogel) $\ll$ $S_\mathrm{BET}$(liogel) $<$ $S_\mathrm{BET}$(aerogel).}
{mild temperature, homogeneous, porosity tunable (aerogels: insulation, Stardust comet-dust collector).}
{shrinkage and cracking, slow gelation; supercritical drying needs an autoclave.}
{oxides, silica shells (Ex.~11b, 13), aerogels (Ex.~11d, 12B).}
\meth{Capping agents \slide{68--73}}
{A surfactant or ligand (e.g.\ a thiol) binds the particle surface during reduction. It limits growth (\textbf{size}), prevents
agglomeration (\textbf{stability}), sets solubility and reactivity (\textbf{surface chemistry}), and, by adsorbing more strongly on some
facets, makes those facets grow slower (\textbf{shape}: rods, cubes, triangles). Thiols (soft bases) bind Au (a soft acid) strongly;
the thiol length tunes the spacing of self-assembled particle arrays. DNA strands link particles into designed assemblies. The shell can
be removed by O$_2$ plasma or calcination.}
{best size and shape control in solution; self-assembly; functional surfaces.}
{the shell blocks the surface (bad for catalysis); removing it can sinter the particles.}
{shaped metals (Ex.~11a, 12A, 16b); AFM capping question (Ex.~17).}
\meth{Synthesis in structured media \slide{74--85}}
{Amphiphiles self-assemble into micelles, which serve as nano-reactors or templates. In a \textbf{microemulsion}, water droplets
inside reverse micelles hold the salt and the reductant (\ce{4Ag+ + BH4- + 4OH- -> 4Ag + BO2- + 4H2O}); the particle cannot grow
larger than the droplet. \textbf{SBA-15}: silica condenses around rod micelles, which are burnt off to leave ordered channels that
host and confine particles. \textbf{Dendrimers} (PAMAM, generations G0, G1\ldots) trap one small particle each (Au@D 1.8 nm,
Au@G1 2.5 nm; larger dendrons resist sintering better). \textbf{Zeolite} cages host exchanged and reduced Ag$^+$.}
{size limited by the cavity; protection against sintering.}
{small amounts; the template must be removed.}
{very small, stable particles.}
\meth{Impregnation \slide{87--88}}
{A porous support (e.g.\ Al$_2$O$_3$ pellets) is soaked in a precursor solution (Ni(NO$_3$)$_2$), dried, calcined to the oxide (NiO)
and reduced in H$_2$ to the metal (Ni).}
{the industrial route to supported catalysts: cheap, simple, any support.}
{broad size distribution; high-temperature steps can sinter the particles.}
{supported catalysts (Ex.~10e, 12B).}
\meth{Galvanic replacement \slide{89--91}}
{A template of a less noble metal is put into a solution of a nobler metal ion. The template is oxidised and dissolves while the noble
metal is reduced onto its surface. This is spontaneous if $E = E^0_\mathrm{reduced} - E^0_\mathrm{oxidised} > 0$. Ag nanowires +
HAuCl$_4$ (\ce{3Ag + Au^{3+} -> 3Ag+ + Au}, $1.50-0.799>0$) give hollow \textbf{Au nanotubes}; Ag nanocubes give nanoboxes and
nanoframes.}
{hollow structures that are hard to make otherwise.}
{needs a suitable pair of reduction potentials.}
{Ex.~15a.}
\meth{Microfluidics \slide{92--100}}
{Pumps inject the precursors into micrometre channels, where an immiscible inert phase (oil or gas) splits the flow into
identical droplets that act as tiny reactors. Separate zones with their own temperature handle nucleation and growth (e.g.\ a core
zone, then a shell zone). The residence time $t = V/Q$ controls growth (CdSe QDs change colour with time, 3 $\to$ 20 min).}
{narrower size distributions than batch mixing; reproducible; low reagent use; continuous.}
{low throughput per channel; design effort and risk of clogging.}
{``reproducible and controlled'' preparations and core--shell design (Ex.~9c, 11b, 13).}

\subsection{Choosing a method}
\begin{center}
\small
\renewcommand{\arraystretch}{1.15}
\begin{tabularx}{\linewidth}{@{}>{\raggedright\arraybackslash}p{5.6cm}X@{}}
\toprule
\textbf{I need\ldots} & \textbf{Use (control parameter)} \\ \midrule
Noble-metal colloid of given size & salt reduction (citrate/reductant ratio, $T$, stirring) \\
Shaped metal (triangles, rods, cubes) & capping agents (facet adsorption) or solvothermal \\
Shaped oxide (cubes, rods) & hydrothermal (NaOH, $T$) \\
Alloy / core--shell & co-reduction / seeded reduction; microreactor for reproducibility \\
Metal@SiO$_2$ & colloid + sol--gel (TEOS) shell, ideally in a microreactor \\
Hollow particles, nanotubes & galvanic replacement \\
Highly porous solid & sol--gel + supercritical drying (aerogel) \\
Ordered macroporous / photonic & polystyrene template + precursor, then burn template \\
Supported catalyst & impregnation $\to$ calcination $\to$ reduction \\
Tonnes of oxide & flame synthesis \\
Ultrathin conformal coating & ALD \\
Carbon nanotubes of given diameter & CVD with catalyst particles of that size \\
\bottomrule
\end{tabularx}
\end{center}
\begin{tipbox}[: answer format for ``propose a method'']
Method + reagents + the parameter that gives the requested feature + one technique to verify it (TEM size/shape, XRD phase,
UV-vis plasmon, XPS oxidation state).
\end{tipbox}
